\section{SF-E836. Complete-intersection point defects with a commuting
generator
triple}\label{SF-E836:sf-e836.-complete-intersection-point-defects-with-a-commuting-generator-triple}

Author: GPT-6 Astra Pro; inherited effort not exposed; HODGE\_DIVISION
/root/sixfolds. Date: 2026-11-21.

\textbf{Status:} complete new proof checked in full by the author,
/root and /root/imprimitive. The root has also read the entire canonical
conversion; the peer audit concerns the complete frozen draft.
All twenty-one display blocks are unchanged by conversion. It extends the reviewed reduced and curvilinear
point-modification obstructions in a precise common
complete-intersection family. No assertion for arbitrary Artinian
modules or all embedding-dimension-four ideals is made. The necessary
companion equations below were previously retained as an unaudited
calculation in \texttt{rm\_artin\_local\_attempts.md}; they are now
proved in full.

\subsection{1. The local statement and its two concrete
corollaries}\label{SF-E836:the-local-statement-and-its-two-concrete-corollaries}

Let \(R=\mathbb C[[x,y,z,w]]\), with maximal ideal \(\mathfrak m\). Let
\(\pi_+,\pi_-\) be two constant rank-two bivectors with complementary
supports. After independent linear choices inside their two planes and
nonzero scalar normalizations, write

\[
 \pi_+=\partial_x\wedge\partial_y,
 \qquad \pi_-=\partial_z\wedge\partial_w.
 \tag{E836.1}
\]

Let \(I\subset R\) be a proper ideal of positive finite colength
generated by a regular sequence \(f_1,f_2,f_3,f_4\). Put \(A=R/I\) and
\(\ell=\dim_{\mathbb C}A>0\). Suppose that

\[
 \{f_j,f_k\}_+\in I\qquad(2\le j<k\le4).
 \tag{E836.2}
\]

Thus a particular triple of the actual ideal generators Poisson-commutes
modulo the ideal for one of the two bivectors. Consider any finite
\(R\)-module with an exact hull-modification sequence

\[
 0\longrightarrow A^n\longrightarrow M\xrightarrow{\phi}R^r
      \longrightarrow A^m\longrightarrow0,
 \qquad m+n>0.
 \tag{E836.3}
\]

Then \(M\) cannot have flat first-order module deformations for both
multiplication classes \(\pi_+\) and \(\pi_-\).

The hypothesis applies in each of the following cases.

\begin{enumerate}
\def\labelenumi{\arabic{enumi}.}
\item
  \textbf{Every common complete-intersection quotient of embedding
  dimension at most three.} The ideal contains a function with a nonzero
  linear part. Its linear part may mix both RM cotangent planes; no
  pure-plane conormal assumption is imposed. Section 5 proves the
  required generator reduction without changing the constant Poisson
  coordinates nonlinearly.
\item
  \textbf{Every presentation with three generators independent of one
  coordinate in one RM plane.} For example, \(I=(x^a,y^b,z^c,w^d)\),
  with arbitrary positive integers \(a,b,c,d\), satisfies (E836.2),
  including the embedding-dimension-four case \(a,b,c,d\ge2\). The first
  generator may depend on all four coordinates; only the indicated three
  need be independent of \(x\).
\end{enumerate}

The second case is a sufficient presentation condition, not a
classification of complete intersections. The remaining
full-embedding-dimension-four complete-intersection gate is not silently
identified with this case.

\subsection{2. Extension matrices and the map to the free deformed
hull}\label{SF-E836:extension-matrices-and-the-map-to-the-free-deformed-hull}

For a bivector \(\pi\), let \(D=\mathbb C[\varepsilon]/(\varepsilon^2)\)
and let \(R_\pi=R\oplus\varepsilon R\) have product

\[
 a*b=ab+\frac{\varepsilon}{2}\pi(da,db).
 \tag{E836.4}
\]

This is the same continuous formal first-order algebra used in E819. In
particular \([a,b]_* =\varepsilon\pi(da,db)\).

If either nonzero bivector deformation exists, E819's
finite-perfect-cone supertrace lemma gives \(\ell m=\ell n\). If either
multiplicity vanishes, this contradicts \(m+n>0\). We therefore assume
\(m=n>0\). The remainder of the proof is independent of the magnitude of
\(m\).

Choose a free-hull basis that identifies its quotient to \(A^m\) with
reduction of the first \(m\) coordinates. Indeed, a surjective matrix
over the local ring \(A\) has a unit \(m\)-minor; lift elementary
operations and that unit minor to \(R\). The image of \(\phi\) is then

\[
 K=I^{\oplus m}\oplus R^{r-m}.
 \tag{E836.5}
\]

The free summand has no \(\operatorname{Ext}^1_R(-,A^m)\), so it splits
in the ordinary extension of \(K\) by \(T=A^m\). No assertion that this
chosen splitting survives a deformation is needed.

Choose elements \(e_i(v)\in M\), depending \(R\)-linearly on
\(v\in R^m\), that lift \(f_i v\) in the first \(m\) hull coordinates.
For every pair there are matrices
\(B_{ij}=-B_{ji}\in\operatorname{Mat}_m(A)\) such that

\[
 f_i e_j(v)-f_j e_i(v)=B_{ij}\bar v\in T.
 \tag{E836.6}
\]

These are exactly the extension matrices: the Koszul resolution of
\(A\), followed by applying \(\operatorname{Hom}_R(-,A^m)\), has all
differentials zero because \(I A^m=0\). Thus

\[
 \operatorname{Ext}^1_R(K,A^m)
   \simeq\operatorname{Ext}^2_R(A^m,A^m)
   \simeq\operatorname{Mat}_m(A)\otimes\bigwedge^2\mathbb C^4.
 \tag{E836.7}
\]

Now assume a flat first-order deformation \(M_\varepsilon\) for \(\pi\).
The map \(\phi\) lifts to an \(R_\pi\)-linear map
\(\phi_\varepsilon:M_\varepsilon\to R_\pi^r\). We recall why the
reviewed E819 lifting applies: finite-length modules over the regular
local fourfold have \(\operatorname{Ext}^j(-,R)=0\) for \(j<4\), by a
finite-length filtration and the four-variable Koszul resolution. The
two exact sequences in (E836.3) therefore give
\(\operatorname{Ext}^1_R(M,R^r)=0\). Lift a finite free resolution and
the degree-zero map. The map-lifting obstruction is precisely this zero
group. Flatness permits the lifted resolution and all kernel
identifications over the dual numbers, as in E819 section 2.1. This uses
actual continuous finite free complexes, not discontinuous operations on
formal power series.

For \(t\in T\), choose a lift \(\widetilde t\in M_\varepsilon\). Its
hull image has the form \(\varepsilon c\), for \(c\in R^r\). Changing
the lift changes \(c\) by an element of \(K\), so its class in \(Q=A^m\)
is well-defined. This defines an \(A\)-linear map

\[
 C_\pi:T\longrightarrow Q,
 \qquad C_\pi\in\operatorname{Mat}_m(A).
 \tag{E836.8}
\]

For linearity, apply the deformed hull map to \(a*\widetilde t\). The
zero reduction of the hull image leaves \(\varepsilon ac\), so the
ordinary scalar action is the induced one. Because \(I\) annihilates
both modules, this factors through \(A\).

\subsection{3. Both companion equations are
necessary}\label{SF-E836:both-companion-equations-are-necessary}

Set

\[
 P^\pi_{ij}=\pi(df_i,df_j)\pmod I\in A.
\]

The previously reviewed hull relation gives

\[
 C_\pi B_{ij}=P^\pi_{ij}I_m.
 \tag{E836.9}
\]

For clarity, lift (E836.6) to the deformed module and apply
\(\phi_\varepsilon\). The hull images of the lifted \(e_i\) are
\(f_i v+\varepsilon a_i\). Their commutator contribution is
\(\varepsilon P^\pi_{ij}v\). The terms \(f_i a_j-f_j a_i\) vanish in the
first \(m\) coordinates modulo \(I\). The correction to the lifted
relation reduces to an element of \(M\), whose hull image is in \(K\),
and vanishes in the same quotient. This proves (E836.9), with the stated
sign convention.

There is also the companion equation in the opposite matrix order:

\[
 B_{ij}C_\pi=P^\pi_{ij}I_m.
 \tag{E836.10}
\]

Here is its full derivation. Since \(f_i t=0\), there exists
\(\widetilde m_i\in M_\varepsilon\) such that
\(f_i*\widetilde t=\varepsilon\widetilde m_i\). Write \(m_i\) for its
reduction. Flatness identifies \(\varepsilon M_\varepsilon\) with \(M\),
so \(m_i\) is determined by this equation. Applying the hull map gives

\[
 \phi(m_i)=f_i c.
 \tag{E836.11}
\]

Split \(c=(c_1,c_2)\) according to the first \(m\) and the remaining
free coordinates. Choose an ordinary splitting of the free part in
(E836.5), denoted \(s:R^{r-m}\to M\). Equation (E836.11) implies

\[
 m_i=e_i(c_1)+s(f_i c_2)+t_i,
 \qquad t_i\in T.
 \tag{E836.12}
\]

Compute the commutator of \(f_i,f_j\) on \(\widetilde t\). On one hand
it is \(\varepsilon P^\pi_{ij}t\). On the other hand it is
\(\varepsilon(f_i m_j-f_j m_i)\). The terms from \(t_i\) vanish because
\(I T=0\), and those from \(s\) cancel by ordinary commutativity. The
remaining terms, by (E836.6), equal \(B_{ij}\bar c_1=B_{ij}C_\pi(t)\).
This proves (E836.10) for every \(t\). No derivative of a chosen lift of
\(c_1\) is omitted: this calculation is already the coefficient of
\(\varepsilon\), where the module action is the ordinary commutative
one.

Apply the two companion equations separately to \(\pi_+\) and \(\pi_-\),
writing \(C_+,C_-\) and \(P^+,P^-\). Associativity gives

\[
 \begin{aligned}
 P^+_{ij}P^-_{kl}I_m
   &=(C_+B_{ij})(C_-B_{kl})\\
   &=C_+(B_{ij}C_-)B_{kl}
     =P^-_{ij}P^+_{kl}I_m.
 \end{aligned}
 \tag{E836.13}
\]

Since \(m>0\), equality of scalar matrices implies the scalar equalities

\[
 P^+_{ij}P^-_{kl}-P^-_{ij}P^+_{kl}=0
          \quad\text{in }A
          \quad\text{for every }i,j,k,l.
 \tag{E836.14}
\]

No coefficient has been divided out in this argument. In particular it
is valid over the nonreduced Artinian ring even when every bracket
coefficient is a zero divisor.

\subsection{4. The commuting triple makes the mixed wedge vanish term by
term}\label{SF-E836:the-commuting-triple-makes-the-mixed-wedge-vanish-term-by-term}

Regard \(P^\pm\) as alternating two-forms on a free \(A\)-module with
basis indexed by the four ideal generators. Hypothesis (E836.2) gives
\(P^+_{jk}=0\) for \(j,k\ge2\). Equation (E836.14), used for the pairs
\((1,i)\) and \((j,k)\) with \(\{i,j,k\}=\{2,3,4\}\), consequently gives

\[
 P^+_{1i}P^-_{jk}=0.
 \tag{E836.15}
\]

Every term of the coefficient of \(P^+\wedge P^-\) is either one of
(E836.15), with a sign, or contains \(P^+_{jk}=0\). Hence that mixed
wedge is zero.

Let \(J=(\partial f_i/\partial z_j)\), with coordinates ordered as
\((x,y,z,w)\). The bracket forms satisfy
\(P^\pm=J\pi_\pm J^{\mathsf t}\) modulo \(I\). Functoriality of the
fourth exterior power gives

\[
 P^+\wedge P^-
    =\det(J)(\pi_+\wedge\pi_-).
 \tag{E836.16}
\]

The right side is nonzero. The two constant bivectors have complementary
supports, and the Jacobian determinant is nonzero in \(A\): for a
zero-dimensional complete intersection over a characteristic-zero field,
the Jacobian ideal is its one-dimensional socle. The exact published
source is Eisenbud--Ulrich, \emph{Duality and socle generators for
residual intersections}, \textbf{J. reine angew. Math. 756 (2019),
183--226, Theorem 7.1, p.210}, with the complete-intersection proof in
\textbf{Appendix A, Theorem A.5, pp.224--225}, using Theorem A.1,
pp.222--223, and Theorem A.3, p.224. The complete invoked statements and
proofs were read in the published author-hosted copy. Theorem 7.1
explicitly uses complete local power-series algebras, so the present
formal setting is included. The quotient is nonzero and finite, and
characteristic zero does not divide its positive length. These are
precisely the hypotheses needed here.

Equations (E836.15)--(E836.16) contradict one another. This proves the
local theorem. Unlike a proposed unrestricted argument from coefficient
minors alone, this proof uses the commuting triple to kill the mixed
wedge term by term; it does not infer proportionality of two forms over
an arbitrary nonreduced ring.

\subsection{5. Every embedding-dimension-at-most-three complete
intersection has the required
presentation}\label{SF-E836:every-embedding-dimension-at-most-three-complete-intersection-has-the-required-presentation}

For \(A=R/I\),
\(\operatorname{edim}A=\dim_{\mathbb C}\mathfrak m/(\mathfrak m^2+I)\).
If this is at most three, there is \(f\in I\) with \(df(0)\ne0\). At
least one of its two RM-plane components is nonzero. Interchange signs
if needed and make a linear change only within that plane so that
\(\partial f/\partial x(0)\ne0\). Rescale the corresponding constant
bivector if necessary. The other plane is unchanged, and both bivectors
still have the form (E836.1). The linear part of \(f\) is allowed to
have components in the other plane.

The element \(f\) is not in \(\mathfrak m I\), since
\(I\subset\mathfrak m\) and \(f\notin\mathfrak m^2\). Extend its nonzero
class to a basis of \(I/\mathfrak m I\), and lift to generators
\(f,h_2,h_3,h_4\) of \(I\). This is an invertible generator change from
any minimal four-generator presentation. These generators are a regular
sequence: in the regular local ring, four generators of the height-four
ideal form a system of parameters and hence a regular sequence.

The formal implicit-function construction gives a unique
\(\varphi\in\mathbb C[[y,z,w]]\), with zero constant term, such that

\[
 f=u\bigl(x-\varphi(y,z,w)\bigr),\qquad u\in R^\times.
 \tag{E836.17}
\]

One may obtain the solution successively by total degree, since the
coefficient of \(x\) in the linear part is nonzero. Substitution makes
\(f\) vanish; formal division by the monic degree-one expression gives
(E836.17), with a unit quotient. In particular \(\varphi\) may have a
nonzero linear \(z\)- or \(w\)-term. We replace the first ideal
generator by \(f_1=x-\varphi\), using the unit \(u\); \textbf{we do not
use \(f_1\) as a new Poisson coordinate}.

For \(j\ge2\), substitute \(x=\varphi\) in \(h_j\), obtaining
\(g_j\in\mathbb C[[y,z,w]]\). The difference \(h_j-g_j\) is divisible by
\(x-\varphi\); subtract this multiple of the first generator. These are
invertible triangular generator operations. We have obtained

\[
 I=(x-\varphi(y,z,w),g_2(y,z,w),g_3(y,z,w),g_4(y,z,w)).
 \tag{E836.18}
\]

Because each \(g_j\) is independent of the unchanged coordinate \(x\),
its bracket with every other \(g_k\) for
\(\pi_+=\partial_x\wedge\partial_y\) is zero as a formal function, not
merely modulo \(I\). Thus (E836.2) holds even for a mixed linear part of
\(f\).

All equations in sections 2--4 may simply be derived for this new
regular sequence. Equivalently their covariance is explicit. For a
generator change \(f'=Uf\), \(U\in\operatorname{GL}_4(R)\),
differentiation modulo \(I\) gives \(df'=Udf\), so
\(P'=UPU^{\mathsf t}\) in \(A\). Taking lifts \(e'_i=\sum_a U_{ia}e_a\)
in the ordinary commutative module gives
\(B'_{ij}=\sum_{a,b}U_{ia}U_{jb}B_{ab}\) with the alternating
convention. This is the second exterior-power change in the Koszul
extension matrices. The map \(C:T\to Q\) is unchanged. Thus both
companion equations are preserved. The Jacobian determinant changes by
\(\det(U)\) modulo \(I\), a unit. No untracked derivative terms survive
modulo \(I\).

\subsection{6. Global balanced consequence and actual noncurvilinear
examples}\label{SF-E836:global-balanced-consequence-and-actual-noncurvilinear-examples}

Let a balanced coherent sheaf on the genuine RM fourfold have a locally
free hull and finite kernel and cokernel of its canonical hull map. At
any point, suppose those two defect modules are free over a common
quotient satisfying the theorem, with positive total multiplicity. The
reviewed E819 characteristic-map argument supplies two global cap-kernel
classes whose bivector components are \(\pi_+,\pi_-\). If the sheaf were
weakly semiregular for a specified coherent effective finite
translation/tensor action, invariant evaluation would vanish on these
classes. Forgetting equivariance and restricting to the completed local
ring gives flat first-order deformations for both bivectors, as proved
there. The local theorem contradicts this. The rational algebraic \(B\)
shear leaves the bivector projections unchanged. No action on a chosen
presentation or on individual local defect bases is assumed.

There are actual noncurvilinear balanced sheaves to which this new
statement applies. Take the rank-32 vector bundle \(H\) constructed and
fully checked in E815/E820, with \(\operatorname{ch}(H)=32v_{1/2}\). At
one point choose a finite subscheme whose completed ideal, in constant
RM-plane coordinates, is

\[
 I_3=(x+z,y^2,z^2,w^2)
       \quad\text{or}\quad
 I_4=(x^2,y^2,z^2,w^2).
 \tag{E836.19}
\]

The first quotient has length eight and embedding dimension three, with
a mixed linear conormal; the second has length sixteen and embedding
dimension four and has the three-generator presentation from case 2.
These are Artinian complete intersections. They define actual algebraic
finite subschemes: each contains a sufficiently high power of the
maximal ideal, and the local ring and its completion have identical
finite quotients. Thus the finitely many required formal jets are
algebraic local functions. No algebraicity of the whole formal logarithm
is asserted.

For either subscheme \(Z\), choose \(1\le m\le32\), a quotient
\(H\twoheadrightarrow Q=\mathcal O_Z^m\), its kernel \(K\), and
\(T=\mathcal O_Z^m\). Such a quotient exists because a vector bundle
over the local Artinian scheme is free. Since \(H\) is locally free and
\(T\) is point-supported, all its positive Ext groups into \(T\) vanish.
The exact sequence for \(K\) and the local Koszul resolution give

\[
 \operatorname{Ext}^1_X(K,T)
   \simeq\operatorname{Ext}^2_X(Q,T),
 \qquad
 \dim_{\mathbb C}\operatorname{Ext}^1_X(K,T)=6m^2\ell>0.
 \tag{E836.20}
\]

Every extension class gives an actual coherent sequence
\(0\to T\to E\to K\to0\). Its torsion is exactly \(T\), because \(K\) is
torsion-free. Its double dual is \(H\), its hull cokernel is \(Q\), and
its entire character is
\(\operatorname{ch}(H)+\operatorname{ch}(T)-\operatorname{ch}(Q)=32v_{1/2}\).
It is ordinarily gluable as a coherent sheaf. The new theorem excludes
weak semiregularity for every specified coherent finite action, if such
an action exists. No simplicity, indecomposability, or action is
asserted.

These examples have equal signed point defect and are not curvilinear.
The first-order obstruction therefore supplies a strict extension of
E819's signed-length statement and E820's common-curvilinear family.
Arbitrary common complete intersections without a commuting triple,
unequal Artinian module types, positive-dimensional singularities,
unstable torsion-free sheaves, and unrelated perfect complexes remain
outside the theorem.

\subsection{7. Source and certificate
boundary}\label{SF-E836:source-and-certificate-boundary}

The new published source is the 2019 Eisenbud--Ulrich paper at the exact
locators above, DOI \texttt{10.1515/crelle-2017-0045}, published primary
PDF:
\url{https://www.math.purdue.edu/~bulrich/papers/published\%20version\%20of\%20Duality\%20for\%20Residual\%20Intersections.pdf}.
Theorem 7.1 and the entire invoked Appendix A proofs were freshly read.
This is an invoked published Jacobian-socle theorem, not an independent
reproof of all upstream different theory.

The finite exact supplement checks the generic twelve-variable
eliminated Jacobian/bracket identity, the noncommutative associative
expansion underlying (E836.13), and the two explicit Artinian socles. It
is not a substitute for the module-lifting, covariance,
formal-generator, global-action, or published-socle arguments. The older
failed matrix solution and the nonintegrable Jacobian counterexample
remain in \texttt{rm\_artin\_local\_attempts.md}; no unrestricted
inference from coefficient minors has been revived.
